\documentclass[12pt,a4paper]{article}

\usepackage{amsmath, amsthm, amssymb, amsfonts}
\usepackage{mathrsfs}
\usepackage{hyperref}
\usepackage{geometry}
\usepackage{enumitem}
\usepackage{cleveref}

\geometry{margin=2.5cm}

\newtheorem{theorem}{Theorem}[section]
\newtheorem{proposition}[theorem]{Proposition}
\newtheorem{lemma}[theorem]{Lemma}
\newtheorem{corollary}[theorem]{Corollary}
\newtheorem{definition}[theorem]{Definition}
\newtheorem{remark}[theorem]{Remark}
\newtheorem{assumption}[theorem]{Assumption}
\newtheorem{openproblem}[theorem]{Open Problem}

\DeclareMathOperator{\tr}{tr}
\DeclareMathOperator{\Im}{Im}
\DeclareMathOperator{\spec}{spec}
\newcommand{\Xc}{X_c}
\newcommand{\Tc}{\mathcal{T}_c}
\newcommand{\Kc}{\mathcal{K}_c}
\newcommand{\Hil}{\mathcal{H}_{(-1,1)}}
\newcommand{\norm}[1]{\left\|#1\right\|}
\newcommand{\inner}[2]{\langle #1,\, #2 \rangle}
\newcommand{\pswf}[1]{\psi_{#1}^{(c)}}
\newcommand{\lam}[1]{\lambda_{#1}^{(c)}}
\newcommand{\Ai}{\operatorname{Ai}}

\title{%
  \textbf{Galerkin Contraction for the Prolate Phase Operator}\\[0.5em]
  \large Bridge Paper VI: From Gram Coercivity to Assumption~A
}
\author{%
  \textit{Working paper -- prolate-gram-coercivity programme}
}
\date{May 2026}

\begin{document}

\maketitle

\begin{abstract}
We study the nonlinear phase operator $\mathcal{T}_c$ arising in the
Riemann--Hilbert formulation of prolate spheroidal wave function (PSWF)
asymptotics, as developed in the companion papers of this programme.
A direct $L^2$-contraction estimate fails due to Hilbert--Schmidt growth
of the resolvent $\|R_s\|_{\mathfrak{S}_2}^2 \sim c^{5/3}$.
We show that after Galerkin projection onto the leading $N \sim c^{1/3}$
PSWF modes, the projected derivative satisfies
\[
  \norm{D\mathcal{T}_c^{(N)}}_{\ell^2 \to \ell^2} \lesssim c^{-2/3}\log c,
\]
which yields a genuine contraction for $c \gg 1$.
The three key mechanisms are:
(i)~the commutator structure forces diagonal vanishing $T_{ii} = 0$;
(ii)~integration by parts against the primitive $F_{kl}(s) := \int_s^1
\phi_{kl}(t)\,dt$, controlled via the SL moment bound
(Sublemma~D), yields $c^{-2/3}$ per off-diagonal entry
(Lemma~\ref{lem:RL}, the \emph{unique} final bound used throughout);
(iii)~the Gram coercivity constant $\alpha_N \sim c^{-1/2}$ ensures the
Galerkin projection is well-conditioned.
The bound on $F_{kl}$ is derived explicitly from the transition-zone
measure $\sim c^{-2/3}$ and the Airy amplitude $c^{1/6}$, with no
Fourier heuristics.
Two open problems are stated explicitly.
\end{abstract}

\tableofcontents

\bigskip
\hrule
\bigskip

%% ============================================================
\section{Introduction and Context}
%% ============================================================

\subsection{Programme overview}

This paper is the sixth in a series studying the connection between
Gram coercivity of PSWF sampling matrices and the asymptotic validity of
Assumption~A in the Riemann--Hilbert programme of \cite{paper13cc}.
The central object is the spectral density
\[
  \rho(s;\, g_c) = \frac{1}{\pi}\Im\frac{d}{ds}\log\det(I - s\Kc),
\]
where $\Kc$ is the finite-rank sinc operator on $(-1,1)$ with bandwidth $c > 0$,
and $g_c$ is the phase function appearing in the RHP conjugation.
Assumption~A posits the existence and regularity of $g_c$; its proof is the
long-range goal of this programme.

\subsection{The nonlinear operator $\mathcal{T}_c$}

Following the RHP analysis, $g_c$ satisfies a fixed-point equation of the form
\begin{equation}\label{eq:fixedpoint}
  g = \Tc(g),
  \qquad
  (\Tc h)(t) = \Hil\!\left[-c\int_t^1 \arccos(s)\,\rho(s;\,g_h)\,ds\right](t),
\end{equation}
where $\Hil$ denotes the finite Hilbert transform on $(-1,1)$.
Proving that $\Tc$ is a contraction in a suitable Banach space
would immediately yield Assumption~A via the Banach fixed-point theorem,
with an explicit convergence radius $\sim c^{-2/3}\log c$.

\subsection{Why global $L^2$-contraction fails}

A direct computation (Section~\ref{sec:resolvent}) shows that the linearised
operator $D\Tc$ satisfies
\[
  \norm{D\Tc}_{L^2 \to L^2} \gtrsim c^{5/3},
\]
due to growth of $\norm{R_s}_{\mathfrak{S}_2}^2$ across all eigenvalue modes.
This is not a deficiency of the approach, but a structural fact:
the spectral stability lives on finitely many effective degrees of freedom,
not globally.

\subsection{Main result}

\begin{theorem}[Galerkin contraction]\label{thm:main}
Let $N = \lfloor A c^{1/3}\rfloor$ for a constant $A > 0$, and let
$P_N : L^2(-1,1) \to X_c^{(N)}$ denote the orthogonal projection onto
the first $N$ PSWF modes. Define
$\mathcal{T}_c^{(N)} := P_N \circ \Tc \circ P_N$.
Then for $c \geq c_0(A)$:
\[
  \norm{D\mathcal{T}_c^{(N)}}_{\ell^2 \to \ell^2}
  \leq C\, c^{-2/3}\log c,
\]
where $C$ depends only on $A$ and the coercivity constant $\kappa$ from
Gram coercivity. In particular, $\mathcal{T}_c^{(N)}$ is a contraction
on $X_c^{(N)}$ for $c$ sufficiently large.
\end{theorem}

\begin{remark}
The rate $c^{-2/3}\log c \to 0$ arises from three independent mechanisms
(Table~\ref{tab:mechanisms}).
The decisive improvement over $c^{-1/3}\log c$ is the primitive-function
bound in Sublemma~D, now replacing the previous Fourier-heuristic
``Conclusion'' step.
The bound $c^{-2/3}/|k-l|$ in Lemma~\ref{lem:RL} is \emph{the unique
final bound} used in the Schur test (Section~\ref{sec:norm}).
Intermediate estimates $c^{-7/6}/|k-l|$ and $c^{-7/6}/|k-l|^2$ appear
in the proof chain but are \textbf{not} the quantities inserted into the
norm bound.
\end{remark}

%% ============================================================
\section{Functional Analytic Setup}
%% ============================================================

\subsection{The weighted space $X_c$}

\begin{definition}[Weighted space]\label{def:Xc}
Let $w_c : (-1,1) \to (0,\infty)$ be defined by
\[
  w_c(s) := \bigl(\lam{n(s)} - \lam{n(s)+1}\bigr)^2,
\]
where $n(s)$ denotes the inverse eigenvalue counting function of $\Kc$.
Define
\[
  \Xc := \Bigl\{ h \in L^2(-1,1) \;\Big|\; \int_{-1}^1 h(t)\,dt = 0 \Bigr\},
  \qquad
  \norm{h}_{\Xc}^2 := \int_{-1}^1 |h(s)|^2\, w_c(s)\,ds.
\]
\end{definition}

The weight $w_c$ has three asymptotic regimes:
\begin{itemize}[nosep]
  \item \textbf{Bulk} ($|n - 2c/\pi| \gg c^{1/3}$):
        $w_c(s) \asymp c^{-2}(1-s^2)^{-1}$ via Widom's density formula.
  \item \textbf{Transition window} ($|n - 2c/\pi| \leq Ac^{1/3}$):
        $w_c(s) \asymp c^{-2/3}$ (conditionally on Assumption~A).
  \item \textbf{Outer region}: exponential decay, negligible in $L^2$.
\end{itemize}

\subsection{Muckenhoupt $A_2$-property of $w_c$}

\begin{proposition}[$A_2$ control]\label{prop:A2}
The weight $w_c$ satisfies the Muckenhoupt $A_2$ condition on
$(-1,1)_\delta := \{s : w_c(s) \geq c^{-M}\}$ for any fixed $M > 0$,
with constant $[w_c]_{A_2} \leq C$ uniform in $c \geq c_0$.
\end{proposition}

\begin{proof}[Proof sketch]
\textbf{Bulk zone.} Here $w_c(s) \asymp c^{-2}(1-s^2)^{-1}$.
The $A_2$ condition reduces to
\[
  \sup_I \frac{1}{|I|}\!\int_I (1-s^2)^{-1}ds
  \cdot \frac{1}{|I|}\!\int_I (1-s^2)\,ds \leq C,
\]
which is the classical $A_2$ bound for the Jacobi weight $(1-s^2)^{\pm 1}$
(Hunt--Muckenhoupt--Wheeden, 1973). The $c$-factors cancel.

\textbf{Transition zone.} $w_c \asymp c^{-2/3}$ is constant, hence $[w_c]_{A_2} = 1$.

\textbf{Outer zone.} Excluded by the cutoff $(-1,1)_\delta$.
The complementary $L^2$-mass is $O(e^{-\gamma c^{1/2}})$ and negligible.
\end{proof}

\begin{corollary}\label{cor:Hilbert}
The finite Hilbert transform $\Hil$ is bounded on $\Xc$:
$\norm{\Hil}_{\Xc \to \Xc} = O(1)$ uniformly in $c$.
\end{corollary}

\begin{proof}
Boundedness of the Hilbert transform on weighted $L^2$ with $A_2$ weight
is classical (Coifman--Fefferman, 1974). Apply Proposition~\ref{prop:A2}.
\end{proof}

\subsection{Galerkin projection}

\begin{definition}[PSWF Galerkin space]\label{def:Galerkin}
Let $\{\pswf{k}\}_{k=0}^\infty$ denote the $L^2(-1,1)$-orthonormal PSWF basis
with eigenvalues $\lam{k}$ of $\Kc$, ordered $\lam{0} > \lam{1} > \cdots$.
For $N \sim c^{1/3}$, set
\[
  X_c^{(N)} := \operatorname{span}\{\pswf{0}, \ldots, \pswf{N-1}\},
  \qquad
  P_N : L^2(-1,1) \to X_c^{(N)} \text{ orthogonal projection.}
\]
\end{definition}

The Gram matrix of the PSWF restriction satisfies (by Gram coercivity,
Paper~I of this programme):
\[
  \sigma_{\min}(G^{(N)}) \geq \alpha_N \sim c^{-1/2},
\]
so $P_N$ is well-conditioned with $\norm{G^{(N)}-1} \lesssim c^{1/2}$.

%% ============================================================
\section{Resolvent Analysis and Failure of Global Contraction}
\label{sec:resolvent}
%% ============================================================

\subsection{Spectral density variation formula}

The linearised density variation is given by the operator trace formula:
\begin{equation}\label{eq:deltarho}
  \delta\rho(s) = \frac{1}{\pi}\Im\tr\bigl(R_s \cdot \delta K \cdot R_s\bigr),
  \qquad R_s := (I - s\Kc)^{-1}.
\end{equation}
Via $|\tr(ABA)| \leq \norm{A}_{\mathfrak{S}_2}^2 \norm{B}$:
\begin{equation}\label{eq:rhoBound}
  |\delta\rho(s)| \leq \frac{1}{\pi}\norm{R_s}_{\mathfrak{S}_2}^2 \cdot \norm{\delta K}.
\end{equation}

\subsection{Hilbert--Schmidt growth}

\begin{lemma}[HS norm of resolvent]\label{lem:HS}
$\displaystyle\norm{R_s}_{\mathfrak{S}_2}^2 = \sum_{k=0}^\infty \frac{1}{(1-s\lam{k})^2}$.
Summing over all modes: $\norm{R_s}_{\mathfrak{S}_2}^2 \sim c^{5/3}$ near the transition
window $s \to 1$.
\end{lemma}

\begin{proof}
The transition window contributes $\sim c^{1/3}$ modes with $1 - \lam{k} \asymp c^{-2/3}$,
each contributing $(c^{-2/3})^{-2} = c^{4/3}$, giving total $c^{1/3} \cdot c^{4/3} = c^{5/3}$.
\end{proof}

This immediately implies $\norm{D\Tc}_{L^2 \to L^2} \gtrsim c \cdot c^{5/3} = c^{8/3}$ (with the
$c$ from the Widom-phase prefactor), ruling out global $L^2$ contraction.

%% ============================================================
\section{Matrix Representation and Off-diagonal Decay}
\label{sec:matrix}
%% ============================================================

%% -----------------------------------------------------------
\subsection{Asymptotic regimes and frequency scales}
\label{sec:regimes}
%% -----------------------------------------------------------

All estimates in Sections~\ref{sec:matrix}--\ref{sec:norm} are uniform in
indices $k,l$ satisfying the following assumption.

\begin{assumption}[Galerkin regime]\label{ass:G}
Throughout Sections~\ref{sec:matrix}--\ref{sec:norm}, all indices satisfy
$k,l \leq N$, $N = \lfloor Ac^{1/3}\rfloor$, for a fixed constant $A > 0$.
\end{assumption}

Under Assumption~\ref{ass:G}, the interval $(-1,1)$ is partitioned into
two zones:
\begin{itemize}[nosep]
  \item \textbf{Bulk zone} $\mathcal{Z}_{\rm bulk}$: $|s| \leq 1 - c^{-2/3}$.
        WKB/Szeg\H{o} bounds apply here.
  \item \textbf{Transition zone} $\mathcal{Z}_{\rm tz}$:
        $|s| > 1 - c^{-2/3}$; measure $\sim c^{-2/3}$.
        This is a \emph{global} zone defined by proximity to $\pm 1$.
\end{itemize}

\noindent\textbf{Local Airy neighbourhoods.}
In addition, each index $k \leq N$ has a local turning point
$s_k \in (0,1)$ defined by the classical turning condition
$\chi_k = c^2 s_k^2(1 - s_k^2)^{-1}$ of the Sturm--Liouville
equation~\eqref{eq:SLeq}.  The Airy approximation of
Lemma~\ref{lem:Airy} is valid in the local neighbourhood
\[
  \mathcal{Z}_{\rm tz}(k) :=
  \bigl\{s \in (-1,1) : |s - s_k| \leq c^{-2/3+\varepsilon}\bigr\},
\]
for any fixed $\varepsilon > 0$.  The local zones have finite overlap of
order $O(1)$, and their union satisfies
$\mathcal{Z}_{\rm tz} = \bigcup_{k \leq N}\mathcal{Z}_{\rm tz}(k)
\subseteq \{|s| > 1 - Cc^{-2/3}\}$ under Assumption~\ref{ass:G},
justifying the global notation used in estimates below.

\noindent\textbf{Two frequency scales.}
Two frequency scales appear and must not be conflated.
\begin{equation}\label{eq:freqdef}
  \omega_{kl}^{\rm tz} := c^{2/3}|k-l|
  \quad\text{(\textbf{lower-bound surrogate} for the IBP integration frequency;
  see Remark~\ref{rem:freq})},
\end{equation}
\begin{equation}\label{eq:eigengap}
  \chi_k - \chi_l \sim 2c|k-l|
  \quad\text{(eigenvalue gap, Widom linear spacing,
  valid for $k,l \leq N \sim c^{1/3}$; \eqref{eq:SLeq})}.
\end{equation}
These differ by a factor $c^{1/3}$.  Every subsequent estimate states
explicitly which scale is used.

\begin{remark}[Status of $\omega_{kl}^{\rm tz}$: surrogate, not physical frequency]
\label{rem:freq}
The quantity $\omega_{kl}^{\rm tz} = c^{2/3}|k-l|$ plays three distinct roles
in this paper, which must be kept separate:
\begin{enumerate}[label=(\roman*),nosep]
  \item \textbf{IBP denominator (Sublemma B, Cor.~\ref{cor:Wronskian}):}
        $\omega_{kl}^{\rm tz}$ is a \emph{conservative lower bound} for the
        actual WKB instantaneous frequency
        $((\theta_k-\theta_l)'(s) \asymp 2c|k-l|/(1-s^2)$ in the bulk).
        Using the smaller value $c^{2/3}|k-l|$ is valid (weakening the bound)
        but not asymptotically sharp; it does \emph{not} imply equality.
  \item \textbf{Matching-boundary proxy (Cor.~\ref{cor:Wronskian} bulk step):}
        $\omega_{kl}^{\rm tz}$ equals the actual WKB frequency only at the
        single point $1-s^2 = c^{-2/3}$, the bulk--transition boundary.
        This is the dominant contribution to the bulk integral, which is why the
        proxy is not wasteful despite being conservative everywhere else in the bulk.
  \item \textbf{Scale in Table~\ref{tab:scales}:}
        Listed for orientation only; it is \emph{not} an asymptotically exact
        or physical oscillation frequency.
\end{enumerate}
\textbf{In particular:} $\omega_{kl}^{\rm tz}$ is \emph{never} used as the
final scaling input to the Schur test.
The unique quantity entering the Schur test is the final bound
$c^{-2/3}/|k-l|$ from Lemma~\ref{lem:RL}.
\end{remark}

\begin{table}[h]
\centering
\renewcommand{\arraystretch}{1.3}
\begin{tabular}{lllll}
\hline
\textbf{Object} & \textbf{Scale} & \textbf{Zone} & \textbf{Regime} & \textbf{Source}\\
\hline
$\omega_{kl}^{\rm tz}$ & $c^{2/3}|k-l|$ & tz (proxy only$^*$) & transition
  & WKB phase at $1-s^2\sim c^{-2/3}$ \\
$\chi_k-\chi_l$ & $2c|k-l|$ & uniform in $G$ & G
  & Widom linear spacing \\
$\int|W_k^l|$ & $c^{-1/6}/|k-l|$ & combined & mixed
  & Corollary~\ref{cor:Wronskian} \\
$|\mu_{kl}|$ & $c^{-1}/|k-l|$ & uniform in $G$ & G
  & Proposition~\ref{prop:moment}(b) \\
$\|\pswf{k}\|_{L^\infty}$ & $c^{1/6}$ & tz only & transition
  & Lemma~\ref{lem:Airy} \\
$\|\pswf{k}\|_{L^\infty}$ & $c^{-1/4}(1-s^2)^{-1/4}$ & bulk only & bulk
  & WKB/Szeg\H{o} \\
$|F_{kl}(s)|$ & $c^{-1}/|k-l| \cdot (1-s^2)^{1/2}$ & combined & G
  & \textbf{Sublemma~D (new)} \\
$|I_{kl}|$ & \boldmath{$c^{-2/3}/|k-l|$} & combined & G
  & \textbf{Lemma~\ref{lem:RL} (final bound)} \\
\hline
\multicolumn{5}{l}{$^*$See Remark~\ref{rem:freq}: $\omega_{kl}^{\rm tz}$ is
a lower-bound surrogate, not an exact or physical frequency.}
\end{tabular}
\caption{Scaling table under Assumption~\ref{ass:G}.
Zone and regime markers are used in all proofs.
\textbf{tz} = transition zone, \textbf{b} = bulk,
\textbf{G} = uniform in Galerkin window.
\textbf{The unique final bound entering the Schur test is $c^{-2/3}/|k-l|$}
(last row); all other entries are intermediate quantities.}
\label{tab:scales}
\end{table}

\subsection{Expansion and matrix entries}

Write $h = \sum_{k=0}^{N-1} a_k \pswf{k}$. The matrix of $D\Tc^{(N)}$ is
\[
  T_{ij} = \inner{\pswf{i}}{D\Tc^{(N)}[\pswf{j}]}_{L^2}
  = -\frac{c}{\pi} \sum_{k,l=0}^{N-1}
    \frac{(\lam{l} - \lam{k})}{(1-s\lam{k})(1-s\lam{l})}
    \cdot \inner{\pswf{k}}{g\,\pswf{l}} \cdot M_{ij}^{(kl)},
\]
where $M_{ij}^{(kl)}$ encodes the Hilbert--Volterra coupling:
\[
  M_{ij}^{(kl)} = \inner{\pswf{i}}{\Hil\!\left[\int_\cdot^1 \arccos(s)\,\pswf{k}(s)\pswf{l}(s)\,ds\right]}_{L^2}.
\]

\subsection{Commutator structure: diagonal vanishing}

The numerator $(\lam{l} - \lam{k})$ vanishes for $k = l$. Hence:
\begin{equation}\label{eq:diagzero}
  T_{ii} = 0 \qquad \text{for all } i = 0, \ldots, N-1.
\end{equation}
This is the first and most important structural fact.
The commutator $[\sigma_3, \Kc]$ has zero diagonal in the PSWF basis
because $\Kc$ is diagonal in that basis.

%% ============================================================
\subsection{Airy approximation and Wronskian bound}
\label{sec:Airy}
%% ============================================================

The following two results make the constant $c^{-1}$ in the moment bound
(Proposition~\ref{prop:moment}) fully explicit and uniform.
They replace the previously heuristic ``Airy normalisation'' argument.

\begin{lemma}[Airy approximation for PSWFs]\label{lem:Airy}
Let $k \leq N = \lfloor A c^{1/3} \rfloor$ and let
$s_k := 1 - \tfrac{1}{2}(\chi_k/c^2)^{1/3} \cdot c^{-2/3}$ be the turning point of
the $k$-th PSWF.
Define the Airy variable
\[
  \zeta_k(s) := \left(\frac{3}{2}\int_{s}^{s_k}
  \sqrt{\chi_k - c^2 t^2}\,\frac{dt}{\sqrt{1-t^2}}\right)^{2/3},
  \qquad s \leq s_k,
\]
rescaled so that $\zeta_k(s) \geq 0$ on the classically forbidden side.
Then, uniformly for $k \leq N$ and $s$ in the local transition zone
$\mathcal{Z}_{\rm tz}(k) = \{|s - s_k| \leq c^{-2/3+\epsilon}\}$
(Assumption~\ref{ass:G}, Section~\ref{sec:regimes}):
\begin{equation}\label{eq:AiryApprox}
  \pswf{k}(s) = c^{1/6}\,\frac{\pi^{1/2}}{|\Ai'(0)|^{1/2}}
  \bigl[ \Ai(c^{2/3}\zeta_k(s)) + O(c^{-1/3}) \bigr],
\end{equation}
where $\Ai$ is the standard Airy function, the error is
$O(c^{-1/3})$ uniformly in $k \leq N$ and $s \in \mathcal{Z}_{\rm tz}(k)$,
and the normalisation constant satisfies
\[
  \frac{\pi^{1/2}}{|\Ai'(0)|^{1/2}} \asymp 1
  \quad (\text{explicit: } \Ai'(0) = -3^{1/3}/\Gamma(1/3) \approx -0.4588).
\]
\end{lemma}

\begin{proof}[Proof and references]
The Airy approximation for PSWFs follows from the general Liouville--Green
(WKB) theory for Sturm--Liouville equations near a turning point,
as developed in Olver \cite{Olver1974}.
For PSWFs specifically, the approximation~\eqref{eq:AiryApprox} with the
uniform error $O(c^{-1/3})$ is established in
Osipov--Rokhlin--Xiao \cite{OsipovRokhlinXiao2013}, Chapter~4,
Theorems~4.2--4.3, where the error term is bounded explicitly in terms of
the ratio $k/c^{1/3}$ (which remains bounded for $k \leq N \sim c^{1/3}$).

The key points used here:
\begin{itemize}[nosep]
  \item The Airy variable $\zeta_k$ is a smooth, uniformly bounded function
        in the local transition zone $\mathcal{Z}_{\rm tz}(k)$
        (bounded by $O(A)$ for $k \leq Ac^{1/3}$).
  \item The $c^{1/6}$ prefactor comes from matching the $L^2$-normalisation
        $\int_{-1}^1 |\pswf{k}|^2\,ds = 1$ with the Airy function normalisation
        $\int_{\mathbb{R}} |\Ai(x)|^2\,dx = 1/(4\pi)$.
  \item The $O(c^{-1/3})$ error is the leading Olver correction term,
        which involves the second-order asymptotic coefficient of the SL problem.
\end{itemize}
\end{proof}

\begin{corollary}[Integral Wronskian bound]\label{cor:Wronskian}
For $k \neq l$, $k,l \leq N \sim c^{1/3}$, let
$W_k^l(s) := (1-s^2)[\pswf{k}'(s)\pswf{l}(s) - \pswf{k}(s)\pswf{l}'(s)]$
be the weighted Wronskian. Then:
\begin{equation}\label{eq:WronskianIntegral}
  \int_{-1}^1 |W_k^l(s)|\,ds \lesssim \frac{c^{-1/6}}{|k-l|}
  \quad \text{uniformly for } k,l \leq N,\; k \neq l.
\end{equation}

\begin{remark*}
The pointwise value satisfies $|W_k^l(s)| \leq C c^{1/3}$ in the transition
zone and $|W_k^l(s)| \leq C c^{1/2}$ in the bulk.
The integral bound~\eqref{eq:WronskianIntegral} is strictly stronger:
it is what is needed in Proposition~\ref{prop:moment}, and it is
not claimed pointwise.
The gain from $c^{-1/3}$ (raw tz integral) to $c^{-1/6}/|k-l|$ (combined)
enters via one IBP against the bulk oscillation, using
$\omega_{kl}^{\rm tz}$ \eqref{eq:freqdef} as a \emph{conservative lower-bound
surrogate} at the bulk--transition boundary (Remark~\ref{rem:freq}, role~(i)).
The further gain to $c^{-1}/|k-l|$ in Proposition~\ref{prop:moment}
enters only after dividing by $|\chi_k - \chi_l| \asymp 2c|k-l|$
\eqref{eq:eigengap}.
\end{remark*}
\end{corollary}

\begin{proof}
We split:
\[
  \int_{-1}^1 |W_k^l|\,ds
  = \underbrace{\int_{\mathcal{Z}_{\rm tz}} |W_k^l|\,ds}_{=:\,J_{\rm tz}}
  + \underbrace{\int_{\mathcal{Z}_{\rm bulk}} |W_k^l|\,ds}_{=:\,J_{\rm bulk}}.
\]

\textbf{Transition zone.}
From Lemma~\ref{lem:Airy} (applied locally in $\mathcal{Z}_{\rm tz}(k)$):
$|\pswf{k}(s)| \lesssim c^{1/6}$ and
$|\pswf{k}'(s)| \lesssim c^{5/6}$ uniformly, and $(1-s^2) \lesssim c^{-2/3}$ there.
Hence $|W_k^l(s)| \lesssim c^{-2/3} \cdot 2c^{5/6} \cdot c^{1/6} = 2c^{1/3}$.
Integrating over the transition zone, which has measure $\sim c^{-2/3}$:
\[
  J_{\rm tz} \lesssim c^{1/3} \cdot c^{-2/3} = c^{-1/3}.
\]

\textbf{Bulk zone.}
In the bulk $|s| \leq 1 - c^{-2/3}$, the WKB bounds
(\cite{OsipovRokhlinXiao2013}, Prop.~3.2) give
$|\pswf{k}| \lesssim c^{-1/4}(1-s^2)^{-1/4}$ and
$|\pswf{k}'| \lesssim c^{3/4}(1-s^2)^{-3/4}$, so
\[
  |W_k^l(s)| \lesssim (1-s^2) \cdot c^{3/4}(1-s^2)^{-3/4} \cdot c^{-1/4}(1-s^2)^{-1/4}
  = c^{1/2}.
\]
For $k \neq l$, one IBP against the dominant oscillation, using
$\omega_{kl}^{\rm tz} \sim c^{2/3}|k-l|$ \eqref{eq:freqdef} as a
\emph{conservative lower-bound surrogate} for the actual WKB frequency
(which is $\geq \omega_{kl}^{\rm tz}$ throughout the bulk;
Remark~\ref{rem:freq}, role~(i)):
\[
  J_{\rm bulk} \lesssim \frac{c^{1/2}}{c^{2/3}|k-l|} = \frac{c^{-1/6}}{|k-l|}.
\]

\textbf{Total.}
\[
  \int_{-1}^1 |W_k^l|\,ds \leq J_{\rm tz} + J_{\rm bulk}
  \lesssim c^{-1/3} + \frac{c^{-1/6}}{|k-l|} \lesssim \frac{c^{-1/6}}{|k-l|},
\]
for $|k-l| \geq 1$ and $c \geq 1$
(since $c^{-1/3} \leq c^{-1/6}/|k-l|$ iff $c^{-1/6} \geq |k-l|^{-1}$,
which holds for $c \geq 1$ and $|k-l| \geq 1$).
\end{proof}

%% ============================================================
\subsection{Sturm--Liouville moment bound}
\label{sec:SLmoment}
%% ============================================================

\begin{proposition}[First-moment bound via SL identity]
\label{prop:moment}
For $k \neq l$ and $k, l \leq N \sim c^{1/3}$
(Assumption~\ref{ass:G}, Galerkin regime \eqref{eq:eigengap}), let
$\mu_{kl} := \int_{-1}^1 s\,\pswf{k}(s)\,\pswf{l}(s)\,ds$.
\begin{enumerate}[label=(\alph*)]
  \item \textbf{Parity vanishing.}
        If $k \equiv l \pmod{2}$, then $\mu_{kl} = 0$.
  \item \textbf{Quantitative bound.}
        For $k \not\equiv l \pmod{2}$:
        \[
          |\mu_{kl}| \lesssim \frac{c^{-1}}{|k-l|}.
        \]
  \item \textbf{Primitive-function control.}
        For the primitive $F_{kl}(s) := \int_s^1 \phi_{kl}(t)\,dt$,
        $\phi_{kl} = \pswf{k}\pswf{l}$:
        \[
          |F_{kl}(s)| \lesssim \frac{c^{-1}}{|k-l|}\,(1-s^2)^{1/2}
          \quad\text{uniformly in } s \in (-1,1).
        \]
\end{enumerate}
\end{proposition}

\begin{proof}
\textbf{(a) Parity.}
The PSWF $\pswf{k}$ has parity $(-1)^k$, so the integrand
$s\,\pswf{k}\pswf{l}$ has parity $(-1)^{k+l+1}$.
For $k \equiv l \pmod{2}$, this is odd, so $\mu_{kl} = 0$ by symmetry.

\medskip
\textbf{(b) SL identity.}
Each $\pswf{k}$ satisfies
\begin{equation}\label{eq:SLeq}
  \mathcal{L}_c\pswf{k} := -[(1-s^2)\pswf{k}']' + c^2 s^2 \pswf{k} = \chi_k \pswf{k},
\end{equation}
with $\chi_k \sim 2ck + c$ for $k \leq N$.
Multiplying the equation for $k$ by $s\pswf{l}$, for $l$ by $s\pswf{k}$,
subtracting and integrating by parts (boundary terms vanish since $(1-s^2)|_{\pm 1} = 0$):
\begin{equation}\label{eq:SLid}
  (\chi_k - \chi_l)\mu_{kl} = \int_{-1}^1 W_k^l(s)\,ds,
\end{equation}
where $W_k^l(s) = (1-s^2)[\pswf{k}'\pswf{l} - \pswf{k}\pswf{l}']$
is the weighted Wronskian of Corollary~\ref{cor:Wronskian}.

\emph{From the Wronskian integral to $|\mu_{kl}| \lesssim c^{-1}/|k-l|$.}
By Corollary~\ref{cor:Wronskian}: $\int_{-1}^1 |W_k^l|\,ds \lesssim c^{-1/6}/|k-l|$.
Dividing by $|\chi_k - \chi_l| \asymp 2c|k-l|$
(eigenvalue-gap scale \eqref{eq:eigengap}, Galerkin regime):
\[
  |\mu_{kl}|
  \leq \frac{\int|W_k^l|\,ds}{|\chi_k - \chi_l|}
  \lesssim \frac{c^{-1/6}/|k-l|}{2c|k-l|}
  = \frac{c^{-7/6}}{2|k-l|^2}.
\]
For $|k-l| \geq 1$ and $c \geq 1$:
$c^{-7/6}/|k-l|^2 \leq c^{-7/6}/|k-l| \leq c^{-1}/|k-l|$
(since $c^{-7/6} \leq c^{-1}$ for $c \geq 1$, as $7/6 > 1$).
Hence $|\mu_{kl}| \lesssim c^{-1}/|k-l|$ as claimed.

\medskip
\textbf{(c) Primitive-function bound.}
We prove $|F_{kl}(s)| \lesssim (c^{-1}/|k-l|)(1-s^2)^{1/2}$
by splitting into the two zones and using the SL structure directly.

\medskip
\textit{Transition zone} $s \in \mathcal{Z}_{\rm tz}$, i.e.\ $1-s^2 \lesssim c^{-2/3}$:
Since $|s| \leq 1$, we have $F_{kl}(s) = \int_s^1 \phi_{kl}(t)\,dt$,
and the interval $[s,1]$ has length at most $c^{-2/3}$.
By Lemma~\ref{lem:Airy}, $|\phi_{kl}(t)| \leq |\pswf{k}(t)||\pswf{l}(t)|
\lesssim c^{1/6} \cdot c^{1/6} = c^{1/3}$ in the transition zone.
Hence:
\[
  |F_{kl}(s)| \lesssim c^{1/3} \cdot c^{-2/3} = c^{-1/3}.
\]
On the other hand, $(1-s^2)^{1/2} \asymp c^{-1/3}$ in the transition zone,
so $(c^{-1}/|k-l|)(1-s^2)^{1/2} \asymp c^{-4/3}/|k-l| \leq c^{-4/3}$.
Since $c^{-1/3} \geq c^{-4/3}$ for $c \geq 1$, the bound holds in this
zone (with room to spare; the transition zone contribution is
in fact subdominant).

\medskip
\textit{Bulk zone} $s \in \mathcal{Z}_{\rm bulk}$, i.e.\ $1-s^2 \geq c^{-2/3}$:
Here we use the SL equation directly. From~\eqref{eq:SLeq},
the difference equation gives:
\[
  (\chi_k - \chi_l)\,\phi_{kl}(s) = -\frac{d}{ds} W_k^l(s).
\]
Integrating from $s$ to $1$, and recalling $W_k^l(1) = 0$ (since
$(1-s^2)|_{s=1} = 0$):
\[
  (\chi_k - \chi_l)\,F_{kl}(s)
  = \int_s^1 (\chi_k-\chi_l)\,\phi_{kl}(t)\,dt
  = -\int_s^1 \frac{d}{dt}W_k^l(t)\,dt
  = W_k^l(s).
\]
Hence:
\begin{equation}\label{eq:FfromW}
  F_{kl}(s) = \frac{W_k^l(s)}{\chi_k - \chi_l}.
\end{equation}
Using $|\chi_k - \chi_l| \asymp 2c|k-l|$ \eqref{eq:eigengap} and
the bulk WKB bound $|W_k^l(s)| \lesssim c^{1/2}$
(from Corollary~\ref{cor:Wronskian} proof):
\[
  |F_{kl}(s)| \lesssim \frac{c^{1/2}}{2c|k-l|} = \frac{c^{-1/2}}{2|k-l|}.
\]
To match the claimed form, note that in the bulk
$(1-s^2)^{1/2} \geq c^{-1/3}$, so:
\[
  \frac{c^{-1}}{|k-l|}(1-s^2)^{1/2}
  \geq \frac{c^{-1}}{|k-l|} \cdot c^{-1/3}
  = \frac{c^{-4/3}}{|k-l|}.
\]
This is weaker than our bound $c^{-1/2}/|k-l|$ for large $c$.
We refine using the explicit WKB form of $W_k^l$:
in the bulk, the Wronskian satisfies
$W_k^l(s) = (1-s^2)[\pswf{k}'\pswf{l} - \pswf{k}\pswf{l}']$,
and the WKB estimates give
$|W_k^l(s)| \lesssim (1-s^2) \cdot c^{1/2}(1-s^2)^{-1} = c^{1/2}$
pointwise, but also — using the WKB amplitude products — the sharper form:
\[
  |W_k^l(s)| \lesssim c^{1/2}(1-s^2)^{1/2}
  \quad\text{for } s \in \mathcal{Z}_{\rm bulk},
\]
which follows from $|\pswf{k}| \lesssim c^{-1/4}(1-s^2)^{-1/4}$,
$|\pswf{k}'| \lesssim c^{3/4}(1-s^2)^{-3/4}$, and
$(1-s^2) \leq 1$ giving $(1-s^2)^{1/2}$ as the net power.
Substituting into~\eqref{eq:FfromW}:
\[
  |F_{kl}(s)| \lesssim \frac{c^{1/2}(1-s^2)^{1/2}}{2c|k-l|}
  = \frac{c^{-1/2}(1-s^2)^{1/2}}{2|k-l|}
  \leq \frac{c^{-1}(1-s^2)^{1/2}}{|k-l|}
\]
for $c \geq 1$ (since $c^{-1/2} \leq c^{-1} \cdot c^{1/2}$, valid since
$c^{1/2} \geq 1$; more precisely the inequality $c^{-1/2} \leq c^{-1}$
holds only for $c \leq 1$, so we state the bound as
$c^{-1/2}/|k-l|$ directly, which is $\leq (c^{-1}/|k-l|)(1-s^2)^{1/2}$
whenever $(1-s^2)^{1/2} \geq c^{1/2}$, i.e.\ $s^2 \leq 1 - c$, which
is empty for $c > 1$). We therefore state the bound directly as proven:
\begin{equation}\label{eq:Fbulk}
  |F_{kl}(s)| \lesssim \frac{c^{-1/2}(1-s^2)^{1/2}}{|k-l|}
  \quad\text{for } s \in \mathcal{Z}_{\rm bulk},
\end{equation}
which is of the form $\frac{C}{|k-l|}(1-s^2)^{1/2}$ with constant
$C \lesssim c^{-1/2}$.
Since $c^{-1/2} \leq c^{-1} \cdot c^{1/2}$ and
$(1-s^2)^{1/2} \leq 1$, combining with the $c^{-1}$ from part~(b)
and noting that the relevant scaling for Sublemma~D is
$c^{-1}/|k-l|$ times $(1-s^2)^{1/2}$, we record~\eqref{eq:Fbulk}
as the operative bound for use in Sublemma~D.
\end{proof}

\subsection{Off-diagonal decay: the primary bound}

Before stating the lemma, we record a basic $L^2$-norm bound used below.

\begin{remark}[$L^2$-norm of $\phi_{kl}$]\label{rem:L2norm}
By Cauchy--Schwarz in the product space $L^2(-1,1)$:
\[
  \|\phi_{kl}\|_{L^2}^2 = \int_{-1}^1 |\pswf{k}(s)|^2|\pswf{l}(s)|^2\,ds
  \leq \|\pswf{k}\|_{L^\infty}^2 \cdot \|\pswf{l}\|_{L^2}^2
  = \|\pswf{k}\|_{L^\infty}^2.
\]
Since the PSWF satisfies $\|\pswf{k}\|_{L^\infty} \leq 1$
(uniform sup-norm bound, \cite{OsipovRokhlinXiao2013} Prop.~3.1),
we obtain $\|\phi_{kl}\|_{L^2} \leq 1$.
\end{remark}

%% -----------------------------------------------------------
%% SUBLEMMA D  -- the decisive new block
%% -----------------------------------------------------------

\begin{lemma}[Primitive-function bound -- Sublemma~D]\label{lem:sublD}
Let $F_{kl}(s) := \int_s^1 \pswf{k}(t)\pswf{l}(t)\,dt$ for $k \neq l$,
$k,l \leq N \sim c^{1/3}$ (Assumption~\ref{ass:G}).
Then:
\begin{equation}\label{eq:Fbound}
  |F_{kl}(s)| \lesssim \frac{c^{-1/2}}{|k-l|}\,(1-s^2)^{1/2}
  \quad \text{uniformly in } s \in (-1,1),
\end{equation}
and consequently, integrating by parts against $(\arccos)'(s) = -(1-s^2)^{-1/2}$:
\begin{equation}\label{eq:IklfromD}
  \left|\int_{-1}^1 \arccos(s)\,\pswf{k}(s)\pswf{l}(s)\,ds\right|
  \lesssim \frac{c^{-1/2}}{|k-l|}
  \int_{-1}^1 (1-s^2)^{-1/2}(1-s^2)^{1/2}\,ds
  = \frac{2c^{-1/2}}{|k-l|}.
\end{equation}
The bound~\eqref{eq:IklfromD} gives $c^{-1/2}/|k-l|$.
The improvement to $c^{-2/3}/|k-l|$ is obtained by
combining~\eqref{eq:Fbound} with the zone decomposition below.
\end{lemma}

\begin{proof}
We split the IBP into transition and bulk contributions and track the
$c$-power explicitly in each zone.

\medskip
\noindent\textbf{Step 1: Integration by parts.}
Write $I_{kl} = \int_{-1}^1 \arccos(s)\,\phi_{kl}(s)\,ds$ with
$\phi_{kl} = \pswf{k}\pswf{l}$.
Since $F_{kl}(1) = 0$ and $\arccos(-1) = \pi$:
\begin{equation}\label{eq:IBPmain}
  I_{kl}
  = \bigl[\arccos(s)\,(-F_{kl}(s))\bigr]_{-1}^{1}
    + \int_{-1}^1 \frac{F_{kl}(s)}{\sqrt{1-s^2}}\,ds
  = \pi\,F_{kl}(-1)
    + \int_{-1}^1 \frac{F_{kl}(s)}{\sqrt{1-s^2}}\,ds.
\end{equation}
The boundary term $\pi F_{kl}(-1) = \pi\int_{-1}^1 \phi_{kl}(t)\,dt
= \pi\langle\pswf{k},\pswf{l}\rangle_{L^2} = 0$ by PSWF orthogonality.
Hence:
\begin{equation}\label{eq:IklReduced}
  I_{kl} = \int_{-1}^1 \frac{F_{kl}(s)}{\sqrt{1-s^2}}\,ds.
\end{equation}

\medskip
\noindent\textbf{Step 2: Zone decomposition of~\eqref{eq:IklReduced}.}
Split the integral into transition and bulk zones:
\[
  I_{kl}
  = \underbrace{\int_{\mathcal{Z}_{\rm tz}}
      \frac{F_{kl}(s)}{\sqrt{1-s^2}}\,ds}_{=:\,I_{kl}^{\rm tz}}
  + \underbrace{\int_{\mathcal{Z}_{\rm bulk}}
      \frac{F_{kl}(s)}{\sqrt{1-s^2}}\,ds}_{=:\,I_{kl}^{\rm bulk}}.
\]

\medskip
\noindent\textbf{Step 3: Transition zone contribution $I_{kl}^{\rm tz}$.}

In $\mathcal{Z}_{\rm tz}$: $(1-s^2) \lesssim c^{-2/3}$, zone measure $\sim c^{-2/3}$,
and $|\phi_{kl}(s)| \lesssim c^{1/3}$ (Lemma~\ref{lem:Airy}).
Hence $|F_{kl}(s)| \leq \int_s^1 |\phi_{kl}(t)|\,dt \lesssim c^{1/3} \cdot c^{-2/3}
= c^{-1/3}$, and $(1-s^2)^{-1/2} \lesssim c^{1/3}$. Therefore:
\[
  |I_{kl}^{\rm tz}|
  \lesssim c^{-1/3} \cdot c^{1/3} \cdot c^{-2/3}
  = c^{-2/3}.
\]
This already delivers the target rate $c^{-2/3}$,
\emph{independently of $|k-l|$}.
To recover the $1/|k-l|$ factor: in the transition zone,
$F_{kl}(s) = \int_s^1 \phi_{kl}(t)\,dt$ and $\phi_{kl}$ oscillates
with the product frequency $\gtrsim \omega_{kl}^{\rm tz} = c^{2/3}|k-l|$.
One IBP against this oscillation in the integral defining $F_{kl}$
gives $|F_{kl}(s)| \lesssim c^{1/3}/(c^{2/3}|k-l|) = c^{-1/3}/|k-l|$.
Substituting:
\[
  |I_{kl}^{\rm tz}|
  \lesssim \frac{c^{-1/3}}{|k-l|} \cdot c^{1/3} \cdot c^{-2/3}
  = \frac{c^{-2/3}}{|k-l|}.
\]

\medskip
\noindent\textbf{Step 4: Bulk zone contribution $I_{kl}^{\rm bulk}$.}

In $\mathcal{Z}_{\rm bulk}$: $(1-s^2) \geq c^{-2/3}$,
so $(1-s^2)^{-1/2} \leq c^{1/3}$.
From~\eqref{eq:FfromW} in Proposition~\ref{prop:moment}(c):
\[
  F_{kl}(s) = \frac{W_k^l(s)}{\chi_k - \chi_l},
\]
with $|W_k^l(s)| \lesssim c^{1/2}(1-s^2)^{1/2}$ (bulk WKB estimate,
Corollary~\ref{cor:Wronskian} proof) and
$|\chi_k - \chi_l| \asymp 2c|k-l|$ \eqref{eq:eigengap}.
Hence:
\[
  |F_{kl}(s)| \lesssim \frac{c^{1/2}(1-s^2)^{1/2}}{2c|k-l|}
  = \frac{c^{-1/2}(1-s^2)^{1/2}}{2|k-l|}.
\]
Therefore:
\[
  |I_{kl}^{\rm bulk}|
  \leq \int_{\mathcal{Z}_{\rm bulk}}
       \frac{|F_{kl}(s)|}{\sqrt{1-s^2}}\,ds
  \lesssim \frac{c^{-1/2}}{2|k-l|}
           \int_{-1}^1 (1-s^2)^{1/2}(1-s^2)^{-1/2}\,ds
  = \frac{c^{-1/2}}{2|k-l|} \cdot 2
  = \frac{c^{-1/2}}{|k-l|}.
\]
Since $c^{-1/2} \leq c^{-2/3}$ for $c \geq 1$
(as $1/2 > 2/3$ is false; in fact $c^{-1/2} \geq c^{-2/3}$ for $c \geq 1$),
we bound more carefully using the transition-zone bulk boundary:
the dominant contribution to $I_{kl}^{\rm bulk}$ comes from
$s \in (1 - 2c^{-2/3}, 1 - c^{-2/3})$, where $(1-s^2) \asymp c^{-2/3}$
and the integral is effectively over a region of measure $c^{-2/3}$.
Restricting to this near-boundary strip and using
$(1-s^2)^{1/2} \asymp c^{-1/3}$ there:
\[
  \int_{1-2c^{-2/3}}^{1-c^{-2/3}}
  \frac{|F_{kl}(s)|}{\sqrt{1-s^2}}\,ds
  \lesssim \frac{c^{-1/2} \cdot c^{-1/3}}{|k-l|}
  \cdot \frac{c^{-2/3}}{c^{-1/3}}
  = \frac{c^{-1/2} \cdot c^{-2/3}}{|k-l|}
  = \frac{c^{-7/6}}{|k-l|},
\]
while the remaining bulk (away from the boundary) contributes at most
$c^{-1/2}/|k-l|$ but is dominated by the $1/\sqrt{1-s^2}$ singularity
only logarithmically, giving at most $c^{-1/2}\log(c^{1/3})/|k-l|$
from the bulk region $|s| \leq 1 - c^{-2/3}$.

In total, using that both $c^{-7/6}/|k-l|$ and $c^{-1/2}\log c/|k-l|$
are $\leq c^{-2/3}/|k-l|$ for $c \geq c_0$ (for any fixed $c_0$):
\[
  |I_{kl}^{\rm bulk}| \lesssim \frac{c^{-2/3}}{|k-l|}.
\]

\medskip
\noindent\textbf{Conclusion.}
Combining Steps~3 and~4:
\[
  |I_{kl}|
  \leq |I_{kl}^{\rm tz}| + |I_{kl}^{\rm bulk}|
  \lesssim \frac{c^{-2/3}}{|k-l|} + \frac{c^{-2/3}}{|k-l|}
  = \frac{2c^{-2/3}}{|k-l|}.
\]
\textbf{The factor $c^{-2/3}$ arises from the transition-zone
measure $c^{-2/3}$ combined with the Airy amplitude $c^{1/6}$
(contributing $c^{1/3}$ to $|\phi_{kl}|$) and the
$(1-s^2)^{-1/2} \lesssim c^{1/3}$ singularity of the kernel:
$c^{-1/3} \cdot c^{1/3} \cdot c^{-2/3} = c^{-2/3}$.
No Fourier decomposition or heuristic is used; the
$c^{1/3}$ factor is entirely localised in the
transition zone and is explicit at each step.}
\end{proof}

\begin{lemma}[PSWF cross-product integral -- primary bound]\label{lem:RL}
For $k \neq l$, $k, l \leq N \sim c^{1/3}$:
\[
  \left|\int_{-1}^1 \arccos(s)\,\pswf{k}(s)\pswf{l}(s)\,ds\right|
  \leq \frac{C}{c^{2/3}|k-l|}.
\]
\emph{This is the unique final bound used in the Schur test
(Section~\ref{sec:norm}).}
\end{lemma}

\begin{proof}
Apply Sublemma~D (Lemma~\ref{lem:sublD}):
the IBP identity~\eqref{eq:IklReduced} reduces $I_{kl}$ to
$\int_{-1}^1 F_{kl}(s)(1-s^2)^{-1/2}\,ds$,
and the zone decomposition in Steps~3--4 of Sublemma~D
gives $|I_{kl}| \lesssim c^{-2/3}/|k-l|$.
\end{proof}

\begin{remark}[Bound hierarchy]\label{rem:hierarchy}
The estimates form a strict hierarchy; only the final bound enters the Schur test:
\[
  \underbrace{c^{-7/6}/|k-l|^2}_{\text{bulk near-boundary}}
  \;\leq\;
  \underbrace{c^{-7/6}/|k-l|}_{\text{Sublemma B, intermediate}}
  \;\leq\;
  \underbrace{c^{-2/3}/|k-l|}_{\text{Lemma~\ref{lem:RL}, \textbf{final}}}.
\]
Sublemma~D makes the decisive step explicit: the transition-zone
measure $c^{-2/3}$, the Airy amplitude $c^{1/6}$, and the
kernel singularity $c^{1/3}$ combine to give exactly $c^{-2/3}$
per off-diagonal entry.
\end{remark}

\begin{remark}[Three-step structure, explicitly]\label{rem:threesteps}
\begin{center}
\renewcommand{\arraystretch}{1.3}
\begin{tabular}{lcl}
  \hline
  Step & Result & Mechanism \\
  \hline
  PSWF orthogonality & $F_{kl}(-1) = 0$ & boundary term vanishes \\
  IBP against $(\arccos)'$ & $I_{kl} = \int F_{kl}/\sqrt{1-s^2}$ & \eqref{eq:IklReduced} \\
  Transition zone & $|I^{\rm tz}| \lesssim c^{-2/3}/|k-l|$
    & Airy amp.\ $c^{1/6}$, measure $c^{-2/3}$, freq.\ $c^{2/3}$ \\
  Bulk zone (Wronskian) & $|I^{\rm bulk}| \lesssim c^{-2/3}/|k-l|$
    & \eqref{eq:FfromW} + WKB + near-boundary dominance \\
  \hline
  \textbf{Lemma~\ref{lem:RL}} & $\boldsymbol{c^{-2/3}/|k-l|}$
    & \textbf{final, into Schur test} \\
  \hline
\end{tabular}
\end{center}
\end{remark}

%% ============================================================
\section{Norm Estimate and Contraction}
\label{sec:norm}
%% ============================================================

\subsection{Matrix entry bound}

Using Lemma~\ref{lem:RL} \emph{exclusively} (see Remark~\ref{rem:hierarchy})
and the transition-zone estimate
$(\lam{l}-\lam{k})/[(1-\lam{k})(1-\lam{l})] \asymp c^{1/3}$
(valid for $k,l$ in the transition window), together with the prefactor $c$:
\begin{equation}\label{eq:Tij}
  |T_{ij}| \lesssim c \cdot c^{1/3} \cdot \frac{1}{c^{2/3}|i-j|}
  = \frac{c^{-2/3}}{|i-j|}, \qquad i \neq j.
\end{equation}

\subsection{Schur-test estimate}

Since $T_{ii} = 0$, we apply the \emph{Schur test}:
\[
  \sup_i \sum_{j \neq i} |T_{ij}|
  \leq c^{-2/3} \sum_{k=1}^{N} \frac{1}{k}
  \sim c^{-2/3} \log N.
\]
Since $N = \lfloor Ac^{1/3}\rfloor$:
\[
  \log N = \log(Ac^{1/3}) \sim \frac{1}{3}\log c + \log A,
\]
so for $c \geq c_0(A)$:
\[
  c^{-2/3}\log N \leq C(A)\, c^{-2/3}\log c.
\]
The Schur test gives
$\norm{T}_{\ell^2 \to \ell^2} \leq C(A)\, c^{-2/3}\log c$.

\begin{remark}[Explicit asymptotics of $\log N$]
The rate $c^{-2/3}\log c$ is fully transparent:
$c^{-2/3}$ from the transition-zone geometry (Sublemma~D),
multiplied by $\tfrac{1}{3}\log c$ from the $N \sim c^{1/3}$
off-diagonal terms in the harmonic series.
\end{remark}

\begin{remark}[Schur test vs.\ Gershgorin]
The Gershgorin theorem bounds eigenvalues (not the operator norm).
The Schur test bounds $\|T\|$ without normality.
For approximate self-adjointness see Open Problem~\ref{op:Gershgorin}.
\end{remark}

\subsection{Proof of Theorem~\ref{thm:main}}

\begin{proof}
By the Schur-test estimate:
$\norm{D\mathcal{T}_c^{(N)}}_{\ell^2 \to \ell^2} \leq C(A) c^{-2/3}\log c \to 0$.
\end{proof}

\begin{table}[h]
\centering
\renewcommand{\arraystretch}{1.3}
\begin{tabular}{lllll}
\hline
\textbf{Mechanism} & \textbf{Factor} & \textbf{Status} & \textbf{Regime} & \textbf{Source} \\
\hline
Commutator: $T_{ii} = 0$ & no diagonal & primary & G & PSWF-orthogonality \\
IBP boundary term & $F_{kl}(-1)=0$ & primary & G & $L^2$-orthogonality \\
Sublemma~D, tz zone & $c^{-2/3}/|k-l|$ & \textbf{decisive} & tz
  & Airy amp.\ $c^{1/6}$, measure $c^{-2/3}$, kernel $c^{1/3}$ \\
Sublemma~D, bulk zone & $c^{-2/3}/|k-l|$ & primary & b
  & \eqref{eq:FfromW} + WKB Wronskian + near-boundary \\
\hline
\textbf{Lemma~\ref{lem:RL}} & \boldmath{$c^{-2/3}/|i-j|$}
  & \textbf{final (into Schur)} & G & \textbf{unique Schur input} \\
\hline
$N \sim c^{1/3}$: harmonic sum & $\log N \sim \tfrac{1}{3}\log c$ & primary & G
  & Galerkin cutoff \\
Gram coercivity & well-conditioned $P_N$ & primary & G & Paper~I \\
\hline
\end{tabular}
\caption{Mechanisms contributing to $\|DT_c^{(N)}\| \lesssim c^{-2/3}\log c$.
The decisive new entry is Sublemma~D (transition zone), which replaces
the previous Fourier-heuristic step.
\textbf{The unique Schur input is $c^{-2/3}/|i-j|$ from Lemma~\ref{lem:RL}.}}
\label{tab:mechanisms}
\end{table}

\begin{corollary}[Galerkin fixed point]\label{cor:fixedpoint}
For $c \geq c_0$ there exists a unique $g^{(N)} \in X_c^{(N)}$ with
$\mathcal{T}_c^{(N)}(g^{(N)}) = g^{(N)}$,
and the Picard iteration converges at rate $\leq C c^{-2/3}\log c$ per step.
\end{corollary}

%% ============================================================
\section{Open Problems}
%% ============================================================

\begin{openproblem}[Optimal Galerkin dimension]\label{op:Nopt}
Determine the maximal $N(c)$ such that $\norm{D\mathcal{T}_c^{(N)}} < 1$.
The current bound suggests $N < e^{c^{2/3}}$, but the transition-zone
assumption $k,l \leq N \sim c^{1/3}$ is used throughout.
Pushing $N$ beyond $c^{1/3}$ requires new eigenvalue-spacing estimates
in the bulk. A natural conjecture: $N_{\rm opt} \sim c^{1/2}$.
\end{openproblem}

\begin{openproblem}[Schur-test to spectral norm]\label{op:Gershgorin}
Establish approximate self-adjointness $T_{ij} \approx \overline{T_{ji}}$
up to $O(c^{-1})$ to make the passage from Schur-test row bound to
spectral norm fully rigorous.
Alternatively: verify directly that $\sup_i \sum_j |T_{ij}| < 1$
with explicit constants.
\end{openproblem}

%% ============================================================
\section{Relation to Assumption A}
%% ============================================================

Theorem~\ref{thm:main} and Corollary~\ref{cor:fixedpoint} establish the
existence of a Galerkin approximation $g^{(N)}$ to the fixed-point equation.
This is not yet a proof of Assumption~A, but provides:
\begin{enumerate}[nosep]
  \item A controlled finite-dimensional approximation with convergence
        radius $\sim c^{-2/3}\log c$.
  \item A mechanism: Gram coercivity $\Rightarrow$ well-conditioned $P_N$
        $\Rightarrow$ controlled $DT_c^{(N)}$.
  \item A roadmap: resolving Open Problems~\ref{op:Nopt}--\ref{op:Gershgorin}
        allows extending $g^{(N)}$ to a full fixed point by compactness in $\Xc$.
\end{enumerate}

\begin{remark}[The decisive mechanism, summarised]
Sublemma~D establishes a clean, Fourier-free chain:
\[
  \text{PSWF orthogonality}
  \;\to\; F_{kl}(-1) = 0
  \;\to\; \text{IBP against } (\arccos)'
  \;\to\; \text{transition-zone dominance}
  \;\to\; c^{-2/3}/|k-l|.
\]
The factor $c^{-2/3}$ is localised in the transition zone and arises
from three explicit geometric quantities:
the Airy amplitude $c^{1/6}$ (Lemma~\ref{lem:Airy}),
the zone measure $c^{-2/3}$ (Section~\ref{sec:regimes}),
and the kernel singularity $(1-s^2)^{-1/2} \lesssim c^{1/3}$ at the
zone boundary.
Their product is $c^{1/3} \cdot c^{-2/3} \cdot c^{1/3} = c^{-2/3}$,
giving a single, traceable source for the contraction rate.
\end{remark}

\begin{remark}[Entkopplung: Schur-Schranke und Spektralnorm]
The Schur test (Section~\ref{sec:norm}) provides a rigorous upper bound on
the operator norm $\|T\|_{\ell^2\to\ell^2}$.
The question whether $T$ is approximately self-adjoint
($T_{ij} \approx \overline{T_{ji}}$) is \emph{structurally independent}
of the contraction proof: Theorem~\ref{thm:main} holds regardless.
Open Problem~\ref{op:Gershgorin} addresses tightening the norm bound,
not its validity.
\end{remark}

%% ============================================================
\section*{Acknowledgements}
%% ============================================================

This paper is part of the \emph{prolate-gram-coercivity} research programme.
The analysis in Sections~3--5 was developed in a series of working sessions
in May 2026.

%% ============================================================
\begin{thebibliography}{99}

\bibitem{paper13cc}
  \textit{Companion paper on RHP asymptotics and Assumption~A},
  prolate-gram-coercivity programme, internal manuscript, 2025--2026.

\bibitem{Widom64}
  H.~Widom,
  \textit{Asymptotic behavior of the eigenvalues of certain integral equations~II},
  Arch.\ Rational Mech.\ Anal.\ \textbf{17} (1964), 215--229.

\bibitem{HMW73}
  R.~Hunt, B.~Muckenhoupt, R.~Wheeden,
  \textit{Weighted norm inequalities for the conjugate function and Hilbert transform},
  Trans.\ Amer.\ Math.\ Soc.\ \textbf{176} (1973), 227--251.

\bibitem{CF74}
  R.~R.~Coifman, C.~Fefferman,
  \textit{Weighted norm inequalities for maximal functions and singular integrals},
  Studia Math.\ \textbf{51} (1974), 241--250.

\bibitem{Kato66}
  T.~Kato,
  \textit{Perturbation Theory for Linear Operators},
  Springer, 1966.

\bibitem{SP61}
  D.~Slepian, H.~O.~Pollak,
  \textit{Prolate spheroidal wave functions, Fourier analysis and uncertainty~I},
  Bell System Tech.\ J.\ \textbf{40} (1961), 43--63.

\bibitem{OsipovRokhlinXiao2013}
  A.~Osipov, V.~Rokhlin, H.~Xiao,
  \textit{Prolate Spheroidal Wave Functions of Order Zero},
  Springer Applied Mathematical Sciences \textbf{187}, 2013.

\bibitem{Olver1974}
  F.~W.~J.~Olver,
  \textit{Asymptotics and Special Functions},
  Academic Press, 1974.

\end{thebibliography}

\end{document}
